% =============================================================================
% APPENDIX CA: LIT-CAMERER: Behavioral Game Theory and Neuroeconomics
% =============================================================================
% Module: BCM2_LIT_CAMERER
% Category: LIT
% Version: 1.0 (January 2026)
% Status: Final
%
% Comprehensive integration of Colin Camerer's research contributions
% on behavioral game theory, cognitive hierarchy, neuroeconomics,
% and strategic decision making into the EBF framework.
%
% Language: English
% =============================================================================

% =============================================================================
% CROSS-REFERENCE STRUCTURE
% =============================================================================
%
% CHAPTER LINKAGE (Required):
% - Primary Chapter: Chapter 9: Complementarity and Interaction (CORE-HOW)
% - Secondary Chapters: Chapter 5: Game-Theoretic Behavior, Chapter 12: Neural Foundations
%
% This appendix integrates with:
%
% DEPENDENCIES (what this appendix REQUIRES):
% - Appendix K (LIT-FEHR): Social preferences foundation
% - Appendix L (LIT-KAHNEMAN): Dual-system framework
% - Appendix B (CORE-HOW): Complementarity parameters
% - Appendix BBB (CORE-WHERE): Parameter repository
%
% DEPENDENTS (what REQUIRES this appendix):
% - Appendix BBB (CORE-WHERE): Parameters for LLMMC
% - All DOMAIN appendices using strategic interaction
% - Appendix AN (METHOD-LLMMC): Estimation procedures
%
% =============================================================================

\begin{tcolorbox}[colback=gray!5!white, colframe=gray!75!black,
    title=Appendix Metadata]
\textbf{Appendix:} CA --- LIT-CAMERER \\
\textbf{Category:} Literature Integration \\
\textbf{Author Focus:} Colin F. Camerer (Caltech) \\
\textbf{Papers:} 22 (95\% Tier-1 Evidence) \\
\textbf{Primary Contribution:} Behavioral game theory, cognitive hierarchy, neuroeconomics
\end{tcolorbox}

\begin{tcolorbox}[colback=purple!5!white, colframe=purple!75!black,
    title=Chapter Linkage]

\textbf{Primary Chapter:} Chapter 9: Complementarity and Interaction
\begin{itemize}[nosep]
\item Section 9.4: Strategic Interaction $\leftrightarrow$ This Appendix Section \ref{sec:ca-bgt}
\item Section 9.6: Bounded Rationality $\leftrightarrow$ This Appendix Section \ref{sec:ca-cognitive}
\end{itemize}

\textbf{Secondary Chapters:}
\begin{itemize}[nosep]
\item Chapter 5: Game-Theoretic Behavior --- cognitive hierarchy model
\item Chapter 12: Neural Foundations --- neuroeconomics integration
\end{itemize}

\textbf{Reading Guide:}
\begin{itemize}[nosep]
\item For conceptual overview: Read Chapter 9 first
\item For formal details: This appendix provides complete specification
\item For proofs: See Appendix D (FORMAL-PROOF)
\end{itemize}

\end{tcolorbox}


\chapter{LIT-CAMERER: Behavioral Game Theory and Neuroeconomics}
\label{app:lit-camerer}


\begin{abstract}
This appendix integrates Colin Camerer's research program into the EBF framework.
Camerer's work establishes that: (1) strategic behavior follows predictable patterns
of bounded rationality through cognitive hierarchy models; (2) economic decisions
have identifiable neural substrates; (3) laboratory games predict field behavior.
His contributions provide the theoretical foundation for modeling strategic
complementarity ($\gamma$) under bounded rationality, linking EBF to neuroscience,
and validating experimental methods. 22 papers spanning 1989-2018 are integrated,
with 95\% providing Tier-1 causal evidence.
\end{abstract}


\begin{tcolorbox}[colback=blue!5!white, colframe=blue!75!black,
    title=Quick Reference: LIT-CAMERER]

\textbf{Core Question:} How do cognitively bounded agents interact strategically, and what neural mechanisms underlie economic decisions?

\textbf{Key Formula (Cognitive Hierarchy):}
\begin{equation}
\sigma_k(a) = BR\left(\sum_{j=0}^{k-1} g_j \cdot \sigma_j\right) \quad \text{where } g_j = \frac{\tau^j e^{-\tau}}{j!}
\end{equation}

\textbf{Key Concepts:}
\begin{itemize}[nosep]
\item \textbf{k-level thinking}: Hierarchical reasoning about others' reasoning
\item \textbf{Neuroeconomics}: Neural substrate of value and choice
\item \textbf{EWA learning}: Experience-weighted attraction in repeated games
\item \textbf{Lab-field validity}: External validity of experimental games
\end{itemize}

\textbf{Cross-References:} Appendix K (LIT-FEHR), L (LIT-KAHNEMAN), B (CORE-HOW), BBB (CORE-WHERE)
\end{tcolorbox}

% -----------------------------------------------------------------------------
% APPENDIX SCOPE BOX
% -----------------------------------------------------------------------------

\begin{tcolorbox}[
    colback=orange!5!white,
    colframe=orange!75!black,
    title={\textbf{Appendix Scope: Ziel / In-Scope / Out-of-Scope / Constraints}},
    fonttitle=\bfseries
]

\textbf{Ziel (Objective):}
\begin{itemize}[nosep]
    \item How do cognitively bounded agents reason about each other's strategies, and what neural mechanisms compute value and guide choice?
\end{itemize}

\textbf{In-Scope:}
\begin{itemize}[nosep]
    \item Theoretical Framework
    \item Axioms and Formal Definitions
    \item Key Results and Findings
    \item Theme 1: Behavioral Game Theory
    \item Theme 2: Cognitive Hierarchy Model
\end{itemize}

\textbf{Out-of-Scope (delegiert an andere Appendices):}
\begin{itemize}[nosep]
    \item LIT-FEHR $\rightarrow$ Appendix K
    \item LIT-KAHNEMAN $\rightarrow$ Appendix L
    \item CORE-HOW $\rightarrow$ Appendix B
    \item CORE-WHERE $\rightarrow$ Appendix BBB
    \item CORE-WHERE $\rightarrow$ Appendix BBB
\end{itemize}

\textbf{Constraints (Anwendungsgrenzen):}
\begin{itemize}[nosep]
    \item [TODO: Define when this appendix does NOT apply]
    \item [TODO: Define required prerequisites]
    \item [TODO: Define known limitations]
\end{itemize}

\textbf{Lieferobjekte (Deliverables):}
\begin{itemize}[nosep]
    \item 6 Table(s)
    \item 16 Equation(s)
    \item 6 Axiom(s)
    \item Worked Example(s)
\end{itemize}

\end{tcolorbox}




\begin{tcolorbox}[colback=red!5!white, colframe=red!75!black,
    title=The Fundamental Question]
\textbf{How do cognitively bounded agents reason about each other's strategies, and what neural mechanisms compute value and guide choice?}
\end{tcolorbox}


%%%%%%%%%%%%%%%%%%%%%%%%%%%%%%%%%%%%%%%%%%%%%%%%%%%%%%%%%%%%%%%%%%%%%%%%%%%%%%%
\section{Theoretical Framework}
\label{sec:ca-theory}
%%%%%%%%%%%%%%%%%%%%%%%%%%%%%%%%%%%%%%%%%%%%%%%%%%%%%%%%%%%%%%%%%%%%%%%%%%%%%%%

Camerer's research program addresses three fundamental gaps in behavioral economics:

\subsection{The Strategic Reasoning Gap}

Standard game theory assumes common knowledge of rationality---all players are rational, know others are rational, know that others know, ad infinitum. Camerer's cognitive hierarchy model shows this assumption fails systematically:
\begin{equation}
\text{Actual reasoning} = \sum_{k=0}^{K} f_k \cdot \text{(k-level reasoning)}
\end{equation}
where most players use $k \leq 2$ levels of reasoning.

\subsection{The Neural Implementation Gap}

Economic models describe choice as if agents compute expected utilities. Camerer's neuroeconomics research asks: what \textit{actually} happens in the brain?
\begin{equation}
V_{\text{neural}}(x) = f_{\text{BOLD}}(\text{vmPFC}, \text{striatum}, \text{insula})
\end{equation}

\subsection{The External Validity Gap}

Laboratory experiments provide control but may lack generalizability. Camerer systematically validates lab-field correspondence:
\begin{equation}
\text{Corr}(\text{Behavior}_{\text{lab}}, \text{Behavior}_{\text{field}}) > 0 \quad \text{(demonstrated)}
\end{equation}


%%%%%%%%%%%%%%%%%%%%%%%%%%%%%%%%%%%%%%%%%%%%%%%%%%%%%%%%%%%%%%%%%%%%%%%%%%%%%%%
\section{Axioms and Formal Definitions}
\label{sec:ca-axioms}
%%%%%%%%%%%%%%%%%%%%%%%%%%%%%%%%%%%%%%%%%%%%%%%%%%%%%%%%%%%%%%%%%%%%%%%%%%%%%%%

Based on Camerer's empirical findings, we formalize six axioms for the EBF:

\begin{tcolorbox}[colback=gray!5!white, colframe=gray!75!black,
    title=Axiom CA-1: Bounded Strategic Depth]
\textbf{Statement:} Human strategic reasoning is bounded to finite levels $k$, with a Poisson distribution of thinking types.
\begin{equation}
f_k = \frac{\tau^k e^{-\tau}}{k!} \quad \text{with } \tau \approx 1.5
\end{equation}
\textbf{Interpretation:} Most people reason 0-2 steps ahead; very few engage in $k > 3$ reasoning.

\textbf{Source:} \citet{PAP-camerer2004cognitive}
\end{tcolorbox}

\begin{tcolorbox}[colback=gray!5!white, colframe=gray!75!black,
    title=Axiom CA-2: Neural Value Representation]
\textbf{Statement:} Subjective value $V(x)$ is encoded in ventromedial prefrontal cortex (vmPFC) activity.
\begin{equation}
\text{BOLD}_{\text{vmPFC}} = \alpha + \beta \cdot V(x) + \epsilon
\end{equation}
\textbf{Interpretation:} Economic value has a common neural currency across reward types.

\textbf{Source:} \citet{PAP-camerer2004neuroeconomics}, \citet{camerer2007brain}
\end{tcolorbox}

\begin{tcolorbox}[colback=gray!5!white, colframe=gray!75!black,
    title=Axiom CA-3: Experience-Weighted Attraction]
\textbf{Statement:} Learning in games follows EWA dynamics combining reinforcement and belief learning.
\begin{equation}
A_i^j(t+1) = \frac{\phi \cdot N(t) \cdot A_i^j(t) + [\delta + (1-\delta) \cdot \mathbb{I}(s_i^t = j)] \cdot \pi_i^j(s_{-i}^t)}{\rho \cdot N(t) + 1}
\end{equation}
\textbf{Interpretation:} People combine past experience with counterfactual reasoning weighted by $\delta$.

\textbf{Source:} \citet{PAP-camerer2005frame}
\end{tcolorbox}

\begin{tcolorbox}[colback=gray!5!white, colframe=gray!75!black,
    title=Axiom CA-4: Overconfidence in Competition]
\textbf{Statement:} Agents overestimate their relative ability, leading to excess market entry.
\begin{equation}
P(\text{entry}) = P(\text{rank} < c) \cdot P(\text{overconfidence}) \quad \text{where } P(\text{overconfidence}) > 0.5
\end{equation}
\textbf{Interpretation:} Overconfidence drives excess competition beyond Nash equilibrium prediction.

\textbf{Source:} \citet{PAP-camerer1999overconfidence}
\end{tcolorbox}

\begin{tcolorbox}[colback=gray!5!white, colframe=gray!75!black,
    title=Axiom CA-5: Theory of Mind Integration]
\textbf{Statement:} Strategic reasoning recruits theory-of-mind neural circuits (TPJ, mPFC).
\begin{equation}
\text{BOLD}_{\text{TPJ/mPFC}} \propto k \quad \text{(level of strategic reasoning)}
\end{equation}
\textbf{Interpretation:} Higher-level strategic thinking requires mentalizing about others' beliefs.

\textbf{Source:} \citet{PAP-camerer2010levels}, \citet{PAP-camerer2017reading}
\end{tcolorbox}

\begin{tcolorbox}[colback=gray!5!white, colframe=gray!75!black,
    title=Axiom CA-6: Lab-Field Correspondence]
\textbf{Statement:} Behavior in controlled experiments correlates with real-world economic decisions.
\begin{equation}
\text{Corr}(\theta_{\text{lab}}, \theta_{\text{field}}) \in [0.3, 0.6] \quad \text{for risk, time, social preferences}
\end{equation}
\textbf{Interpretation:} Laboratory experiments have meaningful external validity for behavioral parameters.

\textbf{Source:} \citet{camerer2016risk}
\end{tcolorbox}


%%%%%%%%%%%%%%%%%%%%%%%%%%%%%%%%%%%%%%%%%%%%%%%%%%%%%%%%%%%%%%%%%%%%%%%%%%%%%%%
\section{Key Results and Findings}
\label{sec:ca-results}
%%%%%%%%%%%%%%%%%%%%%%%%%%%%%%%%%%%%%%%%%%%%%%%%%%%%%%%%%%%%%%%%%%%%%%%%%%%%%%%

\begin{table}[htbp]
\centering
\caption{Summary of Key Results from Camerer's Research}
\label{tab:ca-key-results}
\renewcommand{\arraystretch}{1.3}
\begin{tabular}{@{}llcc@{}}
\toprule
\textbf{Result} & \textbf{Finding} & \textbf{Effect} & \textbf{Tier} \\
\midrule
CA-1 & Strategic depth is bounded ($\tau \approx 1.5$) & k-level distribution & 1 \\
CA-2 & Value encoded in vmPFC & Neural common currency & 1 \\
CA-3 & EWA explains game learning & $\delta \approx 0.5$ & 1 \\
CA-4 & Excess entry from overconfidence & 50\% overentry & 1 \\
CA-5 & Strategic thinking uses ToM circuits & TPJ/mPFC activation & 1 \\
CA-6 & Lab-field correlation exists & $r \in [0.3, 0.6]$ & 1 \\
CA-7 & Trust has neural correlates & Striatum/insula & 1 \\
CA-8 & Expertise reduces bias partially & Domain-specific & 1 \\
\bottomrule
\end{tabular}
\end{table}


%%%%%%%%%%%%%%%%%%%%%%%%%%%%%%%%%%%%%%%%%%%%%%%%%%%%%%%%%%%%%%%%%%%%%%%%%%%%%%%
\section{Theme 1: Behavioral Game Theory}
\label{sec:ca-bgt}
%%%%%%%%%%%%%%%%%%%%%%%%%%%%%%%%%%%%%%%%%%%%%%%%%%%%%%%%%%%%%%%%%%%%%%%%%%%%%%%

\subsection{The Central Insight}

Camerer's magnum opus \textit{Behavioral Game Theory} \citep{PAP-camerer2003behavioral} establishes that people systematically deviate from Nash equilibrium in predictable ways. The key insight is that these deviations are \textit{structured}---not random noise, but systematic patterns that can be modeled.

\subsection{Evidence from Strategic Games}

\paragraph{Beauty Contest Games.} In $p$-beauty contests where players choose numbers hoping to match $p \cdot \text{average}$, Nash equilibrium predicts 0. Actual behavior shows:
\begin{itemize}[nosep]
\item Level-0 players: Choose randomly (mean $\approx 50$)
\item Level-1 players: Best-respond to L0 (choose $p \cdot 50$)
\item Level-2 players: Best-respond to L1 (choose $p^2 \cdot 50$)
\end{itemize}

\paragraph{Ultimatum Games.} Cross-cultural studies \citep{camerer2015ultimatum} show:
\begin{itemize}[nosep]
\item Modal offer: 40-50\% (vs. Nash prediction of minimal offer)
\item Rejection rate for offers $< 20\%$: approximately 50\%
\item Cultural variation exists but basic pattern is universal
\end{itemize}

\subsection{Implications for $\gamma$ (Complementarity)}

For EBF, Camerer's findings imply that strategic complementarity must account for bounded rationality:
\begin{equation}
\gamma_{\text{strategic}}^{\text{bounded}} = \gamma_{\text{Nash}} \cdot \left(1 - \frac{1}{1 + e^{\tau(k-1)}}\right)
\end{equation}

This modifies how we model interaction effects when agents cannot compute full Nash equilibria.


%%%%%%%%%%%%%%%%%%%%%%%%%%%%%%%%%%%%%%%%%%%%%%%%%%%%%%%%%%%%%%%%%%%%%%%%%%%%%%%
\section{Theme 2: Cognitive Hierarchy Model}
\label{sec:ca-cognitive}
%%%%%%%%%%%%%%%%%%%%%%%%%%%%%%%%%%%%%%%%%%%%%%%%%%%%%%%%%%%%%%%%%%%%%%%%%%%%%%%

\subsection{The Model}

The cognitive hierarchy (CH) model \citep{PAP-camerer2004cognitive} provides a tractable alternative to Nash equilibrium:

\begin{definition}[Cognitive Hierarchy Model]
Players are distributed across thinking levels $k \in \{0, 1, 2, ...\}$ with:
\begin{enumerate}[nosep]
\item Level-0 players randomize uniformly
\item Level-$k$ players best-respond to the mixture of levels $0$ to $k-1$
\item The distribution of types follows: $f_k = \frac{\tau^k e^{-\tau}}{k!}$
\end{enumerate}
\end{definition}

\subsection{Key Parameter: $\tau$}

The parameter $\tau$ captures the average depth of reasoning in the population:
\begin{itemize}[nosep]
\item $\tau \approx 1.5$ in most laboratory games
\item This implies most players are Level-1 or Level-2
\item Very few players ($< 5\%$) engage in Level-3+ reasoning
\end{itemize}

\subsection{10C Integration}

The cognitive hierarchy parameter $\tau$ maps to CORE-HOW (complementarity):

\begin{table}[htbp]
\centering
\caption{Cognitive Hierarchy and 10C Framework}
\label{tab:ca-9c}
\renewcommand{\arraystretch}{1.3}
\begin{tabular}{@{}lll@{}}
\toprule
\textbf{10C Question} & \textbf{CH Parameter} & \textbf{Relationship} \\
\midrule
CORE-HOW ($\gamma$) & $\tau$ & Higher $\tau$ → more strategic $\gamma$ \\
CORE-WHEN ($\Psi_C$) & Game complexity & Complex games → lower effective $\tau$ \\
CORE-AWARE ($A$) & Level $k$ & Higher $k$ → higher awareness \\
CORE-WHERE ($\Theta$) & $\tau = 1.5$ (default) & Empirical estimate from experiments \\
\bottomrule
\end{tabular}
\end{table}


%%%%%%%%%%%%%%%%%%%%%%%%%%%%%%%%%%%%%%%%%%%%%%%%%%%%%%%%%%%%%%%%%%%%%%%%%%%%%%%
\section{Theme 3: Neuroeconomics}
\label{sec:ca-neuro}
%%%%%%%%%%%%%%%%%%%%%%%%%%%%%%%%%%%%%%%%%%%%%%%%%%%%%%%%%%%%%%%%%%%%%%%%%%%%%%%

\subsection{The Neuroeconomics Research Program}

Camerer's foundational paper ``Neuroeconomics: Why Economics Needs Brains'' \citep{PAP-camerer2004neuroeconomics} established the field. Key claims:

\begin{enumerate}[nosep]
\item Economic constructs (value, risk, social preference) have neural substrates
\item Brain imaging can inform economic theory
\item Multiple valuation systems compete (Kahneman's System 1/System 2)
\end{enumerate}

\subsection{Neural Value Computation}

The core finding: subjective value is encoded in ventromedial prefrontal cortex (vmPFC):

\begin{equation}
V_{\text{subjective}}(x) \propto \text{BOLD}_{\text{vmPFC}}(x)
\end{equation}

This applies across reward types:
\begin{itemize}[nosep]
\item Money
\item Food
\item Social approval
\item Abstract rewards (fairness)
\end{itemize}

\subsection{Social Neuroeconomics}

Work with Fehr \citep{fehr2007neurobiology} shows social preferences have distinct neural correlates:
\begin{itemize}[nosep]
\item \textbf{Insula}: Activated by unfair offers (disgust response)
\item \textbf{Dorsolateral PFC}: Suppresses fairness concerns when advantageous
\item \textbf{Striatum}: Responds to cooperative outcomes
\end{itemize}


%%%%%%%%%%%%%%%%%%%%%%%%%%%%%%%%%%%%%%%%%%%%%%%%%%%%%%%%%%%%%%%%%%%%%%%%%%%%%%%
\section{Theme 4: Learning in Games}
\label{sec:ca-learning}
%%%%%%%%%%%%%%%%%%%%%%%%%%%%%%%%%%%%%%%%%%%%%%%%%%%%%%%%%%%%%%%%%%%%%%%%%%%%%%%

\subsection{Experience-Weighted Attraction (EWA)}

The EWA model \citep{PAP-camerer2005frame} unifies reinforcement learning and belief learning:

\begin{equation}
A_i^j(t+1) = \frac{\phi \cdot N(t) \cdot A_i^j(t) + [\delta + (1-\delta) \cdot \mathbb{I}(s_i^t = j)] \cdot \pi_i^j(s_{-i}^t)}{\rho \cdot N(t) + 1}
\end{equation}

Key parameters:
\begin{itemize}[nosep]
\item $\phi$: Decay rate of past attractions
\item $\delta$: Weight on foregone payoffs (counterfactual reasoning)
\item $\rho$: Growth rate of experience weight
\end{itemize}

\subsection{Special Cases}

\begin{itemize}[nosep]
\item $\delta = 0$: Pure reinforcement learning
\item $\delta = 1$: Belief learning (fictitious play)
\item Empirical estimate: $\delta \approx 0.5$ (hybrid learning)
\end{itemize}


%%%%%%%%%%%%%%%%%%%%%%%%%%%%%%%%%%%%%%%%%%%%%%%%%%%%%%%%%%%%%%%%%%%%%%%%%%%%%%%
\section{Theme 5: Overconfidence and Market Entry}
\label{sec:ca-overconfidence}
%%%%%%%%%%%%%%%%%%%%%%%%%%%%%%%%%%%%%%%%%%%%%%%%%%%%%%%%%%%%%%%%%%%%%%%%%%%%%%%

\subsection{The Phenomenon}

Camerer and Lovallo \citep{PAP-camerer1999overconfidence} demonstrate that overconfidence causes excess market entry:

\begin{enumerate}[nosep]
\item When entry depends on \textit{skill}, people overenter relative to Nash equilibrium
\item When entry depends on \textit{luck}, entry rates match Nash
\item The difference reveals overconfidence about relative ability
\end{enumerate}

\subsection{Quantitative Findings}

\begin{itemize}[nosep]
\item Excess entry: $\approx 50\%$ above Nash prediction
\item Overconfidence about rank: Better-than-average effect
\item Persists with experience and incentives
\end{itemize}

\subsection{Implications for DOMAIN Applications}

For EBF domain applications, overconfidence must be modeled when:
\begin{itemize}[nosep]
\item Market entry decisions (entrepreneurship)
\item Competition for scarce resources
\item Performance-contingent pay
\end{itemize}


%%%%%%%%%%%%%%%%%%%%%%%%%%%%%%%%%%%%%%%%%%%%%%%%%%%%%%%%%%%%%%%%%%%%%%%%%%%%%%%
\section{Integration with Other Appendices}
\label{sec:ca-integration}
%%%%%%%%%%%%%%%%%%%%%%%%%%%%%%%%%%%%%%%%%%%%%%%%%%%%%%%%%%%%%%%%%%%%%%%%%%%%%%%

\subsection{Connection to Appendix K (LIT-FEHR)}

Camerer and Fehr have extensively collaborated, particularly on social neuroeconomics:

\begin{tcolorbox}[colback=yellow!5!white, colframe=yellow!75!black,
    title=Integration: LIT-CAMERER $\leftrightarrow$ LIT-FEHR]
\textbf{What LIT-CAMERER provides:}
\begin{itemize}[nosep]
\item Neural mechanisms for social preferences
\item Cognitive hierarchy for strategic settings
\item Laboratory validation methodology
\end{itemize}

\textbf{What LIT-FEHR provides:}
\begin{itemize}[nosep]
\item Social preference foundations (inequity aversion, reciprocity)
\item Field experiment paradigms
\item Institutional design applications
\end{itemize}
\end{tcolorbox}

\subsection{Connection to Appendix L (LIT-KAHNEMAN)}

Camerer builds on Kahneman's dual-system framework:

\begin{itemize}[nosep]
\item System 1: Intuitive, fast → Level-0 and Level-1 reasoning
\item System 2: Deliberative, slow → Level-2+ reasoning
\item Neural implementation: Limbic vs. prefrontal systems
\end{itemize}

\subsection{Connection to Appendix B (CORE-HOW)}

\begin{equation}
\gamma_{\text{Camerer}} = \gamma_{\text{Nash}} \cdot g(\tau) \quad \text{where } g(\tau) \text{ captures bounded rationality}
\end{equation}


%%%%%%%%%%%%%%%%%%%%%%%%%%%%%%%%%%%%%%%%%%%%%%%%%%%%%%%%%%%%%%%%%%%%%%%%%%%%%%%
\section{Worked Example: Predicting Game Behavior}
\label{sec:ca-example}
%%%%%%%%%%%%%%%%%%%%%%%%%%%%%%%%%%%%%%%%%%%%%%%%%%%%%%%%%%%%%%%%%%%%%%%%%%%%%%%

\subsection{Setup}
Consider a firm designing a competitive bidding mechanism. Using Camerer's cognitive hierarchy model to predict bid distributions.

\paragraph{The Game.} First-price sealed-bid auction with $n = 4$ bidders, values uniformly distributed on $[0, 100]$.

\paragraph{Nash Prediction.} Risk-neutral Nash equilibrium bid: $b^*(v) = \frac{n-1}{n} v = 0.75v$.

\paragraph{Application with Cognitive Hierarchy.}

\textbf{Step 1: Define Level-0 behavior.}
Level-0 bidders bid randomly: $b_0 \sim U[0, 100]$.

\textbf{Step 2: Compute Level-1 response.}
Level-1 bidders best-respond to Level-0:
\begin{equation}
b_1(v) = \frac{1}{2} \cdot \mathbb{E}[b_0] = \frac{1}{2} \cdot 50 = 25 \text{ (for all } v > 25\text{)}
\end{equation}

\textbf{Step 3: Compute Level-2 response.}
Level-2 bidders best-respond to mixture of L0 and L1:
\begin{equation}
b_2(v) = \text{BR}\left(\frac{\tau^0 e^{-\tau}}{0!} \cdot b_0 + \frac{\tau^1 e^{-\tau}}{1!} \cdot b_1\right)
\end{equation}

\textbf{Step 4: Aggregate prediction.}
With $\tau = 1.5$:
\begin{itemize}[nosep]
\item 22\% are Level-0
\item 33\% are Level-1
\item 25\% are Level-2
\item 20\% are Level-3+
\end{itemize}

\paragraph{Result.}
The cognitive hierarchy model predicts:
\begin{itemize}[nosep]
\item Average bids will be \textit{lower} than Nash (due to L0 and L1 players)
\item Bid distribution will be \textit{bimodal} (L0 randomizers + strategic types)
\item Revenue will be 10-15\% below Nash prediction
\end{itemize}

\paragraph{Discussion.}
This example shows how EBF can use cognitive hierarchy to improve predictions in strategic settings, informing mechanism design with realistic assumptions about reasoning.


%%%%%%%%%%%%%%%%%%%%%%%%%%%%%%%%%%%%%%%%%%%%%%%%%%%%%%%%%%%%%%%%%%%%%%%%%%%%%%%
\section{Implications for Practice}
\label{sec:ca-practice}
%%%%%%%%%%%%%%%%%%%%%%%%%%%%%%%%%%%%%%%%%%%%%%%%%%%%%%%%%%%%%%%%%%%%%%%%%%%%%%%

\begin{tcolorbox}[colback=green!5!white, colframe=green!75!black,
    title=Practical Implications]
\begin{enumerate}
\item \textbf{Mechanism Design:} Design mechanisms robust to bounded rationality (assume $\tau \approx 1.5$)
\item \textbf{Market Entry:} Expect 50\% excess entry when success depends on relative skill
\item \textbf{Learning Interventions:} Use EWA parameters to predict adaptation speed
\item \textbf{Neural Targeting:} Consider neural mechanisms when designing financial products
\item \textbf{Lab Experiments:} Trust lab results for preference parameters ($r \approx 0.3-0.6$ with field)
\end{enumerate}
\end{tcolorbox}


%%%%%%%%%%%%%%%%%%%%%%%%%%%%%%%%%%%%%%%%%%%%%%%%%%%%%%%%%%%%%%%%%%%%%%%%%%%%%%%
\section{Parameters for LLMMC Estimation}
\label{sec:ca-llmmc}
%%%%%%%%%%%%%%%%%%%%%%%%%%%%%%%%%%%%%%%%%%%%%%%%%%%%%%%%%%%%%%%%%%%%%%%%%%%%%%%

From Camerer's research, the following parameters are available for LLMMC:

\begin{table}[htbp]
\centering
\caption{LLMMC Parameters from LIT-CAMERER}
\label{tab:ca-llmmc}
\renewcommand{\arraystretch}{1.3}
\begin{tabular}{@{}llccc@{}}
\toprule
\textbf{Parameter} & \textbf{Description} & \textbf{Point Estimate} & \textbf{95\% CI} & \textbf{Tier} \\
\midrule
$\tau$ & Cognitive hierarchy depth & 1.5 & [1.0, 2.0] & 1 \\
$\delta_{\text{EWA}}$ & Foregone payoff weight & 0.5 & [0.3, 0.7] & 1 \\
$P(\text{overentry})$ & Excess entry rate & 0.50 & [0.40, 0.60] & 1 \\
$r_{\text{lab-field}}$ & Lab-field correlation & 0.40 & [0.30, 0.60] & 1 \\
$k_{\text{modal}}$ & Modal reasoning level & 1 & [1, 2] & 1 \\
\bottomrule
\end{tabular}
\end{table}


%%%%%%%%%%%%%%%%%%%%%%%%%%%%%%%%%%%%%%%%%%%%%%%%%%%%%%%%%%%%%%%%%%%%%%%%%%%%%%%
\section{Summary}
\label{sec:ca-summary}
%%%%%%%%%%%%%%%%%%%%%%%%%%%%%%%%%%%%%%%%%%%%%%%%%%%%%%%%%%%%%%%%%%%%%%%%%%%%%%%

\begin{tcolorbox}[colback=green!10!white, colframe=green!75!black,
    title=\textbf{Appendix CA: Summary}]

\textbf{Core Question:} How do bounded agents reason strategically and what are the neural mechanisms?

\textbf{Main Contributions:}
\begin{itemize}[nosep]
    \item Cognitive hierarchy model ($\tau \approx 1.5$) for strategic interaction
    \item Neuroeconomics: neural substrates of value (vmPFC) and social cognition (TPJ/mPFC)
    \item EWA learning model unifying reinforcement and belief learning
    \item Overconfidence causes 50\% excess market entry
    \item Lab-field correlation validates experimental methods
\end{itemize}

\textbf{Key Outputs:}
\begin{itemize}[nosep]
    \item $\tau$: Average depth of strategic reasoning ($\approx 1.5$)
    \item $\delta$: EWA foregone payoff weight ($\approx 0.5$)
    \item Neural value: $V \propto \text{BOLD}_{\text{vmPFC}}$
\end{itemize}

\textbf{Integration:} This appendix connects to CORE-HOW (strategic complementarity),
CORE-WHEN (cognitive load context), and LIT-FEHR/LIT-KAHNEMAN via shared neural foundations.

\end{tcolorbox}


%%%%%%%%%%%%%%%%%%%%%%%%%%%%%%%%%%%%%%%%%%%%%%%%%%%%%%%%%%%%%%%%%%%%%%%%%%%%%%%
\section{Glossary of Symbols}
\label{sec:ca-glossary}
%%%%%%%%%%%%%%%%%%%%%%%%%%%%%%%%%%%%%%%%%%%%%%%%%%%%%%%%%%%%%%%%%%%%%%%%%%%%%%%

\begin{tcolorbox}[colback=blue!5!white, colframe=blue!75!black]
\textbf{Central Reference:} For complete symbol definitions, see
\textbf{Appendix G} (\textit{Glossary and Notation}).

This section lists only symbols introduced or specialized in this appendix.
\end{tcolorbox}

\vspace{0.5em}

\begin{table}[htbp]
\centering
\caption{Symbols Introduced in Appendix CA}
\label{tab:ca-symbols}
\renewcommand{\arraystretch}{1.3}
\begin{tabular}{@{}clp{6cm}c@{}}
\toprule
\textbf{Symbol} & \textbf{Name} & \textbf{Definition} & \textbf{Also in G?} \\
\midrule
$\tau$ & CH depth parameter & Mean of Poisson distribution over thinking levels & NEW \\
$k$ & Thinking level & Number of steps of strategic reasoning & NEW \\
$f_k$ & Level distribution & Proportion of population at level $k$ & NEW \\
$\delta$ & EWA weight & Weight on foregone payoffs in learning & NEW \\
$\phi$ & EWA decay & Decay rate for past attractions & NEW \\
$\gamma_{\text{strategic}}$ & Strategic complementarity & Complementarity in strategic settings & \checkmark \\
\bottomrule
\end{tabular}
\end{table}


%%%%%%%%%%%%%%%%%%%%%%%%%%%%%%%%%%%%%%%%%%%%%%%%%%%%%%%%%%%%%%%%%%%%%%%%%%%%%%%
\section{Critical Foundations and Objections}
\label{sec:ca-foundations}
%%%%%%%%%%%%%%%%%%%%%%%%%%%%%%%%%%%%%%%%%%%%%%%%%%%%%%%%%%%%%%%%%%%%%%%%%%%%%%%

\subsection{Objection 1: Neuroeconomics is Neuroimaging, Not Economics}

\textbf{Objection:} Brain scans don't add predictive value beyond behavioral measures. ``Seeing the brain light up'' doesn't improve economic models.

\textbf{Response:} Camerer's neuroeconomics program addresses this through:
\begin{itemize}[nosep]
\item Neural data resolves model ambiguity when behavioral data can't distinguish theories
\item Brain measures predict behavior \textit{before} it occurs (anticipatory signals)
\item Understanding mechanisms enables better interventions
\end{itemize}

\subsection{Objection 2: Cognitive Hierarchy is Arbitrary}

\textbf{Objection:} The Poisson distribution for thinking levels is assumed, not derived. Why not normal, exponential, or other distributions?

\textbf{Response:}
\begin{itemize}[nosep]
\item Poisson fits data better than alternatives (AIC comparisons)
\item Parameter $\tau$ is stable across games ($\approx 1.5$)
\item The model's predictive success justifies the distributional assumption
\end{itemize}

\subsection{Objection 3: Lab Results Don't Generalize}

\textbf{Objection:} Laboratory games with students and small stakes don't predict real economic behavior.

\textbf{Response:} Camerer explicitly tests and demonstrates:
\begin{itemize}[nosep]
\item Lab-field correlations of 0.3-0.6 for risk and social preferences
\item Replication across cultures (ultimatum games worldwide)
\item Stakes effects are small in most domains
\end{itemize}


%%%%%%%%%%%%%%%%%%%%%%%%%%%%%%%%%%%%%%%%%%%%%%%%%%%%%%%%%%%%%%%%%%%%%%%%%%%%%%%
\section{Open Issues and Research Agenda}
\label{sec:ca-open}
%%%%%%%%%%%%%%%%%%%%%%%%%%%%%%%%%%%%%%%%%%%%%%%%%%%%%%%%%%%%%%%%%%%%%%%%%%%%%%%

\begin{table}[htbp]
\centering
\caption{Research Roadmap for Camerer-Related Questions}
\label{tab:ca-roadmap}
\renewcommand{\arraystretch}{1.3}
\begin{tabular}{@{}clcl@{}}
\toprule
\textbf{\#} & \textbf{Issue} & \textbf{Priority} & \textbf{Status} \\
\midrule
1 & Context-dependence of $\tau$ & HIGH & Partial evidence \\
2 & Neural-behavioral prediction gap & HIGH & Active research \\
3 & Learning speed heterogeneity & MEDIUM & Underexplored \\
4 & Overconfidence interventions & MEDIUM & Mixed results \\
5 & Cross-cultural $\tau$ estimates & LOW & Limited data \\
\bottomrule
\end{tabular}
\end{table}


%%%%%%%%%%%%%%%%%%%%%%%%%%%%%%%%%%%%%%%%%%%%%%%%%%%%%%%%%%%%%%%%%%%%%%%%%%%%%%%
\section{Complete Paper Inventory}
\label{sec:ca-papers}
%%%%%%%%%%%%%%%%%%%%%%%%%%%%%%%%%%%%%%%%%%%%%%%%%%%%%%%%%%%%%%%%%%%%%%%%%%%%%%%

\begin{table}[htbp]
\centering
\caption{All 22 Camerer Papers in EBF Database}
\label{tab:ca-papers}
\renewcommand{\arraystretch}{1.2}
\scriptsize
\begin{tabular}{@{}llll@{}}
\toprule
\textbf{ID} & \textbf{Year} & \textbf{Theme} & \textbf{Tier} \\
\midrule
camerer1989empirical & 1989 & Behavioral Finance & 1 \\
camerer1997neuroeconomics & 1997 & Game Theory & 1 \\
camerer1999overconfidence & 1999 & Overconfidence & 1 \\
camerer2003behavioral & 2003 & Behavioral Game Theory & 1 \\
camerer2004neuroeconomics & 2004 & Neuroeconomics & 1 \\
camerer2004cognitive & 2004 & Cognitive Hierarchy & 2 \\
camerer2005frame & 2005 & EWA Learning & 1 \\
camerer2006hyperbolic & 2006 & Time Preferences & 1 \\
camerer2007brain & 2007 & Neuroeconomics & 1 \\
fehr2007neurobiology & 2007 & Social Neuroeconomics & 1 \\
camerer2008advice & 2008 & Advice Taking & 1 \\
camerer2009strategic & 2009 & Strategic Thinking & 1 \\
camerer2010levels & 2010 & Theory of Mind & 1 \\
loewenstein2010visceral & 2010 & Hedonic Experience & 1 \\
camerer2011trust & 2011 & Trust Games & 1 \\
camerer2012expertise & 2012 & Expert Judgment & 1 \\
camerer2013social & 2013 & Neuroeconomics & 1 \\
camerer2014prediction & 2014 & Prediction Markets & 1 \\
camerer2015ultimatum & 2015 & Cross-Cultural Games & 1 \\
camerer2016risk & 2016 & Risk Preferences & 1 \\
camerer2017reading & 2017 & Mind Reading & 1 \\
camerer2018prospect & 2018 & Prospect Theory & 1 \\
\bottomrule
\end{tabular}
\end{table}


%%%%%%%%%%%%%%%%%%%%%%%%%%%%%%%%%%%%%%%%%%%%%%%%%%%%%%%%%%%%%%%%%%%%%%%%%%%%%%%
\section{References}
\label{sec:ca-references}
%%%%%%%%%%%%%%%%%%%%%%%%%%%%%%%%%%%%%%%%%%%%%%%%%%%%%%%%%%%%%%%%%%%%%%%%%%%%%%%

\begin{tcolorbox}[colback=blue!5!white, colframe=blue!75!black]
\textbf{Master Bibliography:} For complete reference details, see
\texttt{bcm\_master.bib} in the repository.

This section lists key references for this appendix.
\end{tcolorbox}

\vspace{0.5em}

\subsection{Key References}

\begin{itemize}[nosep]
    \item \textbf{Behavioral Game Theory:} \citet{PAP-camerer2003behavioral}
    \item \textbf{Cognitive Hierarchy:} \citet{PAP-camerer2004cognitive}
    \item \textbf{Neuroeconomics:} \citet{PAP-camerer2004neuroeconomics}, \citet{fehr2007neurobiology}
    \item \textbf{EWA Learning:} \citet{PAP-camerer2005frame}
    \item \textbf{Overconfidence:} \citet{PAP-camerer1999overconfidence}
\end{itemize}

% Link to master bibliography
\nocite{bcm_master}


%%%%%%%%%%%%%%%%%%%%%%%%%%%%%%%%%%%%%%%%%%%%%%%%%%%%%%%%%%%%%%%%%%%%%%%%%%%%%%%
% CROSS-REFERENCE FOOTER (Required)
%%%%%%%%%%%%%%%%%%%%%%%%%%%%%%%%%%%%%%%%%%%%%%%%%%%%%%%%%%%%%%%%%%%%%%%%%%%%%%%

\vspace{1em}
\hrule
\vspace{0.5em}

\noindent\textit{Cross-references:}
Appendix G (Central Glossary),
Appendix K (LIT-FEHR),
Appendix L (LIT-KAHNEMAN),
Appendix B (CORE-HOW),
Appendix BBB (CORE-WHERE).

\noindent\textit{Master Bibliography:} \texttt{bcm\_master.bib}


% =============================================================================
% END OF APPENDIX CA
% =============================================================================
